\documentclass[11pt,a4paper]{amsart}
\usepackage[T1]{fontenc}
\usepackage{amssymb}
\usepackage[margin=1in]{geometry}
\usepackage[hidelinks]{hyperref}
\hypersetup{pdftitle={No fully polynomial additive scheme for calibration distance unless P = NP}}
\urlstyle{same}
\DeclareFontShape{OT1}{cmr}{bx}{n}{<-> cmbx10}{}
\DeclareMathSizes{8}{8}{8}{8}
\DeclareMathSizes{9}{9}{8}{8}
\DeclareMathSizes{10.95}{10.95}{8}{8}
\newtheorem{theorem}{Theorem}
\newcommand{\caldist}{\operatorname{CalDist}}
\title{No fully polynomial additive scheme for calibration distance unless P = NP}
\date{}
\begin{document}
\begin{abstract}
For an explicit rational distribution on a finite feature space, the
distance from a supplied predictor to the nearest exactly calibrated
deterministic predictor has a fully polynomial additive approximation
scheme if and only if \(\mathrm P=\mathrm{NP}\). The hardness holds
with noiseless binary labels and nonuniform masses. A reduction from
strongly NP-complete 3-PARTITION forces heavy anchor pairs into distinct
blocks. Each block has the correct calibrated mean at its target
integer load. Any imbalance creates a normalized inverse-polynomial gap.
\end{abstract}
\maketitle

\section{The approximation problem}

Let \(D\) be a rational distribution on \(\mathcal X\times\{0,1\}\),
where \(\mathcal X\) is finite. We may discard zero-mass atoms.
Write \(w_x=\Pr(X=x)\), \(\mu_x=\mathbb E[Y\mid X=x]\),
and let \(f:\mathcal X\to[0,1]\) be the given rational forecast.
We seek
\[
 \caldist_D(f)=
 \min_{g:\,\mathbb E[Y\mid g(X)]=g(X)}
       \sum_{x\in\mathcal X}w_x\,|f(x)-g(x)|.
\]
The repaired predictor \(g\) is a deterministic function of \(x\).
Every feature atom is assigned whole to one output. The repair may
send two atoms with the same original forecast to different outputs.

By the partition characterization of Qiao \cite[Fact~1]{qiao},
each partition implements a calibrated predictor by assigning to
a block \(S\) its mass-weighted label mean
\[
 p_S=\frac{\sum_{x\in S}w_x\mu_x}{\sum_{x\in S}w_x}.
\]
If two blocks have the same mean, their union has the same mean and
their contributions to the objective are unchanged. Conversely,
the level sets of a calibrated predictor form a partition of this
kind. Thus the minimum is attained and can be studied by checking
partitions.

Qiao \cite[Theorems~1--2]{qiao} proves an additive PTAS,
an exact polynomial-time algorithm when the masses are
uniform and labels noiseless, and exact NP-hardness when either
restriction is removed. His Section~4 asks whether there is an
additive FPTAS, as recorded in catalog Q460 \cite{catalog}.
The exact hardness uses subset sum and supplies no
inverse-polynomial approximation gap on a strongly hard family.

\begin{theorem}\label{main}
There is a deterministic algorithm which, for all such rational
inputs and rational \(\epsilon>0\), gives the distance to additive error
\(\epsilon\) in time polynomial in the input bit length and
\(1/\epsilon\) if and only if \(\mathrm P=\mathrm{NP}\).
The hardness direction uses only noiseless labels.
\end{theorem}

\section{The reduction}

Take a strongly NP-complete 3-PARTITION instance with \(3m\) positive
integers \(a_i\), target \(B\), total \(mB\), and
\(B/4<a_i<B/2\) \cite{gj}. We use unnormalized masses and divide
the resulting loss by the total mass. Make one distinct ordinary
atom of weight \(a_i\), label \(1\) and forecast \(0\) for each item.
For \(j=1,\ldots,m\) use the parameters
\[
 q_j=\frac14+\frac{j}{12m},\qquad
 M=120m^2B,\qquad
 \mu_j=q_j-\frac{(1-q_j)B}{M}.
\]
For each \(j\) create a pair of anchor atoms, both forecast \(q_j\).
Give the label-\(1\) anchor weight \(M\mu_j\) and the label-\(0\)
anchor weight \(M(1-\mu_j)\). The bounds \(1/4\le q_j\le1/3\)
and \(B/M\le1/120\) give \(\mu_j>1/5\) and \(1-\mu_j\ge2/3\).
The normalization is \(T=m(M+B)\). Let \(K=B\sum_jq_j\).

Consider a partition with unnormalized loss below \(K+2/3\).
If one block contains anchors of different indices, its loss
on those two atoms is at least
\[
   (M/5)|q_j-q_k|\ \ge\ 2mB
\]
by the triangle inequality. This exceeds \(K+2/3\), using
\(K\le mB/3\) and \(mB\ge1\). If there is no such block but a
pair is split, its positive-label anchor belongs to a block of
label-\(1\) atoms alone. That block's prediction is \(1\), so
the anchor contributes at least
\((M/5)(1-q_j)\ge 2M/15>K+2/3\).
Consequently a partition below the cutoff keeps every pair
intact and keeps different pairs apart. It may have ordinary-only
blocks in addition to the \(m\) anchored blocks.

\section{The loss from an incorrect load}

Let \(s_j\) be the ordinary mass in anchor block \(j\),
and let \(s_0\) be the mass outside all anchored blocks.
An intact pair with ordinary mass \(s\) has mean
and unnormalized loss
\[
 p_j(s)=\frac{M\mu_j+s}{M+s},\qquad
 F_j(s)=s\,p_j(s)+M|p_j(s)-q_j|.
\]
Here
\(p_j(s)-q_j=(1-q_j)(s-B)/(M+s)\).
At \(s=B\) the loss is \(Bq_j\).
For \(s\ge B\), substituting the positive sign gives
\(F_j(s)=Bq_j+(s-B)\).
For \(s\le B\), substitution gives
\[
 F_j(s)
   =s q_j+\frac{(1-q_j)(B-s)(M-s)}{M+s}
   \ \ge\ Bq_j-\tfrac13(B-s).
\]
The inequality uses \(M>B\) and \(q_j\le 1/3\).

Write \(L_-=\sum_j(B-s_j)_+\) and
\(L_+=\sum_j(s_j-B)_+\), with \(x_+=\max\{x,0\}\).
Since \(\sum_js_j+s_0=mB\), we have \(L_-=L_++s_0\).
An ordinary-only block has prediction \(1\) and loss equal
to its mass. Thus a partition with all pairs intact and
separate has total loss at least
\[
 K-\tfrac13 L_-+L_++s_0=K+\tfrac23 L_-.
\]
The loads are sums of indivisible integer weights.
If \(L_-=0\), then \(s_0=0\), \(L_+=0\), and every load is
\(B\). These assignments form a 3-partition because the size
bounds force exactly three items in each block. Such a partition
has loss \(K\). If there is no 3-partition, then \(L_-\ge1\).
Together with the costs of mixed or split anchors this proves
\[
 \begin{array}{ll}
 \text{YES:}&\caldist_D(f)=K/T,\\[2pt]
 \text{NO:}&\caldist_D(f)\ge (K+2/3)/T.
 \end{array}
\]

\section{Bit complexity and converse}

All emitted masses and forecasts are rational with polynomial
bit length. In the strongly hard restriction of 3-PARTITION,
\(B\) is polynomially bounded, and \(T=m(120m^2B+B)\) is too.
Run the putative additive FPTAS with
\(\epsilon=1/(6T)\). A value within that error lies on the
correct side of \((K+1/3)/T\), hence decides those instances
in polynomial time. The same argument would put \(\mathrm{NP}\)
in \(\mathrm{BPP}\) if the approximator used bounded-error
randomization.

The rational threshold decision belongs to \(\mathrm{NP}\):
name a partition and evaluate its block means and costs exactly.
Polynomial bit length follows by summing and multiplying the
explicit rational numerators and denominators for finitely many
atoms. If \(\mathrm P=\mathrm{NP}\), binary search on
thresholds in \([0,1]\) gives the required error after
\(O(\log(1/\epsilon))\) comparisons for \(0<\epsilon<1\).
A witness can also be recovered by asking whether a prefix
of the partition's certificate has an extension at the chosen cutoff.
This proves Theorem~\ref{main}.

The theorem concerns explicit weighted atoms. The unit-weight
uniform/noiseless special case stays polynomial by \cite{qiao}.
The proof leaves the uniform/noisy approximation question outside its scope.
Fractional distance and repairs of split unit identities
\cite{qz,wy} use different feasible sets. The result is
consistent with an additive PTAS and does not address
sample-based estimation or tolerant testing.

\section*{Review and disclosure}
Internally reviewed, not peer reviewed. No proof-assistant
formalization is included. Under human direction, ChatGPT Astra 6.2
(OpenAI) generated the discovery proof and original replay.
OpenAI Codex performed the independent mathematical review.
GPT-6.2 Sol (OpenAI) derived the derivative-free loss bound,
checked the literature and prepared the revised manuscript,
expanded replay and worked illustration. DannyExperiments is
the project's citation identity. The manuscript has no byline.

\begin{thebibliography}{9}
\bibitem{qiao} M. Qiao, Computational and statistical hardness of
calibration distance, 2026.
\href{https://arxiv.org/abs/2603.18391v2}{arXiv:2603.18391v2}.
\bibitem{catalog} Open Problems Catalog, Q460.
\href{https://github.com/10m-u/open-problems-catalog/blob/ca8d4d44ba3db527f19f28928c240168f545ecfe/problems/computer_science_ai/classical-statistical-machine-learning.md\#q460}{Pinned entry}.
\bibitem{gj} M. R. Garey and D. S. Johnson, Strong NP-completeness
results: motivation, examples, and implications,
\emph{J. ACM} 25 (1978), 499--508.
\href{https://doi.org/10.1145/322077.322090}{DOI}.
\bibitem{qz} M. Qiao and L. Zheng, On the distance from calibration in
sequential prediction, \emph{Proc. COLT}, PMLR 247 (2024).
\href{https://proceedings.mlr.press/v247/qiao24a.html}{Proceedings}.
\bibitem{wy} Z. Wang and X. Yang, Sharp integrality gaps in calibration
distance, 2026.
\href{https://arxiv.org/abs/2610.05679v1}{arXiv:2610.05679v1}.
\end{thebibliography}
\end{document}
